\documentclass[10pt,a4paper,oneside]{article}
\usepackage[utf8]{inputenc}
\usepackage[english,russian]{babel}
\usepackage{amsmath}
\usepackage{amsthm}
\usepackage{amssymb}
\usepackage{enumerate}
\usepackage{stmaryrd}
\usepackage{comment}
\usepackage{cmll}
\usepackage{mathrsfs}
\usepackage{hyperref}
\usepackage[left=2cm,right=2cm,top=2cm,bottom=2cm,bindingoffset=0cm]{geometry}
\usepackage{proof}
\usepackage{tikz}
\usepackage{multicol}

\makeatletter
\newcommand{\dotminus}{\mathbin{\text{\@dotminus}}}

\newcommand{\@dotminus}{%
  \ooalign{\hidewidth\raise1ex\hbox{.}\hidewidth\cr$\m@th-$\cr}%
}
\makeatother

\usetikzlibrary{arrows,backgrounds,patterns,matrix,shapes,fit,calc,shadows,plotmarks}

\newtheorem{definition}{Определение}
\begin{document}

\begin{center}{\Large\textsc{\textbf{Теоретические домашние задания}}}\\
             \it Теория типов, ИТМО, М3232-М3239, весна 2026 года\end{center}

\section*{Домашнее задание №1: лямбда исчисление --- бестиповое и просто-типизированное}


Для проверки и демонстрации заданий используйте какой-нибудь эмулятор лямбда-исчисления,
например LCI: \url{https://www.chatzi.org/lci/}

\begin{enumerate}

\item Постройте лямбда выражения:
\begin{enumerate}
\item деление чёрчевского нумерала на 3.
\item ограниченное вычитание 1 (указание: воспользуйтесь упорядоченными парами).
\item целочисленное деление.
\end{enumerate}

\item На лекции был приведён комбинатор неподвижной точки $Y := \lambda f.(\lambda x.f\ (x\ x))\ (\lambda x.f\ (x\ x))$, обладающий свойством
$Y\ P =_\beta P\ (Y\ P)$ для любого терма $P$. С его помощью оказывается возможным реализовывать рекурсию.

Например, зададим функцию, возводящую 2 в соответствующую степень: $$P := \lambda f.\lambda x.(IsZero\ x)\ 1\ ((f\ (Dec\ x)) \cdot 2)$$
Сравните это с кодом на Си: 
\begin{center}\verb!unsigned f (unsigned x) { return x == 0 ? 1 : f (x-1) * 2; }!\end{center}
Тогда, вызванная как $Y\ P\ x$, эта функция вычислит $2^x$. Например, $Y\ P\ 1 =_\beta$

\vspace{-0.3cm}
$$\begin{array}{l}=_\beta P\ (Y\ P)\ 1 = (\lambda f.\lambda x.(IsZero\ x)\ 1\ ((f (Dec\ x))\cdot 2)))\ (Y\ P)\ 1 \\
  =_\beta (IsZero\ 1)\ 1\ ((Y\ P\ (Dec\ 1)) \cdot 2))) =_\beta (Y\ P\ 0) \cdot 2 \\
  =_\beta (P\ (Y\ P)\ 0) \cdot 2 \\
  =_\beta (IsZero\ 0)\ 1\ ((Y\ P\ (Dec\ 0)) \cdot 2))) \cdot 2\\
  =_\beta 1 \cdot 2 =_\beta 2\end{array}$$

С помощью $Y$-комбинатора реализуйте:
\begin{enumerate}
\item Вычисление $k$-го простого числа.
\item Частичный логарифм.
\item Предложите три других комбинатора неподвижной точки (других --- то есть, не бета-экви\-ва\-лент\-ных $Y$ и между собой).
\end{enumerate}

\item Напомним, что список может быть задан с помощью алгебраического типа с двумя конструкторами, \verb!Nil! и \verb!Cons!
(см. доказательство неразрешимости исчисления предикатов). С учётом этого знания, и с учётом представления 
алгебраических типов, приведённого на лекции, реализуйте следующие конструкции:
\begin{enumerate}
\item Функцию, вычисляющую длину списка.
\item Функцию высшего порядка \verb!map2! --- последовательно применяет функцию к головам двух списков, возвращая список результатов:
\verb!map2 (*) [1;3] [2;4]! вернёт \verb![2;12]!.
\item Функцию \verb!rev!, возвращающую перевёрнутый список. Например, $\text{rev}[1,3,5] = [5,3,1]$.
\end{enumerate}


\item Покажите, что для чёрчевских нумералов $\overline{\lambda a.\lambda b.b\ a} =_\beta \overline{a^b}$.

\item Бесконечное количество комбинаторов неподвижной точки. Дадим следующие определения
$$\begin{array}{l}
L := \lambda abcdefghijklmnopqstuvwxyzr.r(thisisafixedpointcombinator)\\
R := LLLLLLLLLLLLLLLLLLLLLLLLLL\end{array}$$
В данном определении терм $R$ является комбинатором неподвижной точки: каков бы ни был терм
$F$, выполнено $R\ F =_\beta F\ (R\ F)$.
\begin{enumerate}
\item Докажите, что данный комбинатор --- действительно комбинатор неподвижной точки.
\item Пусть в качестве имён переменных разрешены русские буквы. Постройте аналогичное выражение
по-русски: с 32 параметрами и осмысленной русской фразой в терме $L$; покажите, что оно является
комбинатором неподвижной точки.
\end{enumerate}

%\item Найдите необитаемый тип в просто-типизированном лямбда-исчислении (напомним: тип $\tau$ 
%называется необитаемым, если ни для какого $P$ не выполнено $\vdash P : \tau$).
\item (задача пропущена)

\begin{comment}
\item Напомним определение: комбинатор --- лямбда-выражение без свободных переменных. Также напомним:
$$\begin{array}{l}
S := \lambda x.\lambda y.\lambda z.x\ z\ (y\ z)\\
K := \lambda x.\lambda y.x\\
I := \lambda x.x
\end{array}$$

Известна теорема о том, что для любого комбинатора $X$ можно найти выражение $P$
(состоящее только из скобок, пробелов и комбинаторов $S$ и $K$), что $X =_\beta P$.
Будем говорить, что комбинатор $P$ \emph{выражает} комбинатор $X$ в базисе $SK$.

Косвенным аргументом (пояснением, но не доказательством!) в пользу этой теоремы являются 
два следующих соображения:
\begin{itemize}
\item теорема о замкнутости ИфИИВ: если $\vdash \varphi$, то $\vdash_\rightarrow \varphi$,
значит, если выражение имеет тип, то этот тип можно получить с помощью доказательства в
стиле Гильберта;
\item типы комбинаторов $S$ и $K$ --- это, соответственно, вторая и первая схемы аксиом.
\end{itemize}

Докажите тип следующих выражений как логическое высказывание с помощью
гильбертового вывода и, пользуясь этим доказательством как источником вдохновения, 
выразите комбинаторы в базисе $SK$:
\begin{enumerate}
\item $\lambda x.\lambda y.\lambda z.y$
\item $\lambda x.\lambda y.\lambda z.y x z$
\item $\overline{1}$
\item $Not$
\item $Xor$
\item $InR$
\end{enumerate}
\end{comment}

\item Сформулируем определение бета-редукции на языке пред-лямбда-термов. $A \rightarrow_\beta B$, если:
\begin{itemize}
\item $A \equiv (\lambda x.P)\ Q$, $B \equiv P\ [x := Q]$, при условии свободы для подстановки;
\item $A \equiv (P\ Q)$, $B \equiv (P'\ Q')$, при этом $P \rightarrow_\beta P'$ и $Q = Q'$, либо $P = P'$ и $Q \rightarrow_\beta Q'$;
\item $A \equiv (\lambda x.P)$, $B \equiv (\lambda x.P')$, и $P \rightarrow_\beta P'$.
\end{itemize}

\begin{enumerate}
\item Найдите лямбда-выражение, бета-редукция которого не может быть произведена из-за нарушения
правила свободы для подстановки (для продолжения редукции потребуется производить переименование
связанных переменных). Поясните, какое ожидаемое ценное свойство будет нарушено, если ограничение
правила проигнорировать.
\item Покажите, что недостаточно наложить требования на исходное выражение, и свобода для подстановки
может быть нарушена уже в процессе редукции исходно полностью корректного лямбда-выражения.
\end{enumerate}

\begin{comment}
\item Будем говорить, что выражение $A$ находится в \emph{слабой заголовочной нормальной форме} (WHNF),
если оно не имеет вид $A \equiv (\lambda x.P)\ Q$ (то есть, самый верхний терм его не является редексом).
Выражение находится в \emph{заголовочной нормальной форме} (HNF), когда его верхний терм --- не редекс и не лямбда-абстракция
с бета-редексами в теле.
\begin{enumerate}
\item Приведите нормальным порядком редукции выражение $2\ 2$ в СЗНФ.
\item Приведите нормальным порядком редукции выражение $Y\ (\lambda f.\lambda x.(IsZero\ x)\ 1\ (x \cdot f(x-1)))\ 3$ в СЗНФ.
\item Верно ли, что <<нормальность>> формы выражения может в процессе редукции только усиливаться
(никакая --- слабая заголовочная Н.Ф. --- заголовочная Н.Ф. --- нормальная форма)?
\end{enumerate}
\end{comment}

%\item Две задачи на вычисление СЗНФ при помощи нормального порядка редукции.
%
%\begin{enumerate}
%\item Постройте функцию прибавления 1 к значениям из списка в лямбда-исчислении:\\ \verb!let plus1 l = map (fun x -> x+1) l!.
%Постройте бесконечный список из нечётных чисел \verb![1,3,..]!. Примените функцию \verb!plus1! к списку и получите результат
%в СЗНФ.
%
%\item Напишите функцию вычисления суммы первых двух элементов списка:\\
%\verb!let compute (l1::l2::ls) = (l1+l2, ls)!, примените её к результату предыдущего пункта,
%получите результат в СЗНФ.
%\end{enumerate}

\item Как мы уже разбирали, $\not\vdash x\ x:\tau$ в силу дополнительных ограничений
правила
$$\infer[x \notin FV(\Gamma)]{\Gamma, x:\tau\vdash x:\tau}{}$$

Найдите лямбда-выражение $N$, что $\not\vdash N:\tau$ в силу ограничений правила
$$\infer[x \notin FV(\Gamma)]{\Gamma \vdash \lambda x.N:\sigma\to\tau}{\Gamma, x:\sigma \vdash N:\tau}$$

%\item Верно ли, что $S = B(BW)(BBC)$? Если нет, то как правильно?

\item Рассмотрим подробнее отличия исчисления по Чёрчу и по Карри. 
Определим точно бета-редукцию в исчислении по Чёрчу: $A \rightarrow_\beta B$, если

\begin{tabular}{ll}
$A = (\lambda x^\sigma.P)\ Q$, & $B = P [x := Q]$\\
$A = P\ Q$, & $B = P\ Q'$ или $B = P'\ Q$ при $P \rightarrow_\beta P'$ и $Q \rightarrow_\beta Q'$\\
$A = \lambda x^\sigma.P$, & $B = \lambda x^\sigma.P'$ при $P \rightarrow_\beta P'$
\end{tabular}

\begin{enumerate}
\item Покажите, что в исчислении по Карри не выполняется свойство расширения типизации
(subject expansion) даже в такой формулировке: 
если $\vdash_\text{к} M:\sigma$, $M \twoheadrightarrow_\beta N$ и $\vdash_\text{к} N:\tau$,
то необязательно, что $\sigma=\tau$.
\item Покажите, что в исчислении по Чёрчу свойство расширения типизации в такой формулировке также не выполняется:

\begin{center}найдутся такие $M,N,\sigma$, что $\vdash_\text{ч} N:\sigma$, $M\twoheadrightarrow_\beta N$, но $\not\vdash_\text{ч} M:\sigma$.\end{center}

Но при этом в исчислении по Чёрчу выполнено свойство расширения типизации в такой формулировке:

\begin{center}если $\vdash_\text{к} M:\sigma$, $M \twoheadrightarrow_\beta N$ и $\vdash_\text{к} N:\tau$,
то тогда $\sigma=\tau$.\end{center}
\end{enumerate}
\end{enumerate}

\section*{Задание №2. Просто-типизированное лямбда-исчисление}

\begin{enumerate}
\item Докажите (по возможности строго), что предложенный на лекции алгоритм преобразования лябмда-выражения
в выражение в базисе SKI действительно заканчивает работу, строя выражение, содержащее только комбинаторы
SKI и скобки. В частности, ожидается, что вы полностью укажете схему индуктивного доказательства
(в известном автору доказательстве две индукции с содержательными условиями --- одна вложена в другую).

Остальные части доказательства теоремы (например, что полученное выражение 
бета-эквивалентно исходному) доказывать не обязательно, если они очевидны, т.е. вы понимаете,
что нужно доказать, какова схема доказательства и как в целом доказывать отдельные переходы.

\item Покажите, что следующая система комбинаторов образует полный базис в бестиповом
лямбда-исчислении, но соответствующая им система аксиом в исчислении гильбертового типа
не образует полного базиса для импликативного фрагмента:

$$\begin{array}{l}
I := \lambda x.x\\
K := \lambda x.\lambda y.x\\
S' := \lambda i.\lambda x.\lambda y.\lambda z.i\ (i\ ((x\ (i\ z))\ (i\ (y\ (i\ z)))))
\end{array}$$

Указание: покажите невыводимость $(\varphi \rightarrow \varphi \rightarrow \psi) \rightarrow (\varphi \rightarrow \psi)$.

\item Не пользуясь алгоритмом (предлагаемое алгоритмом выражение скорее всего всё равно будет громоздким), 
выразите следующие лямбда-термы в базисе $SKI$:
\begin{enumerate}
\item $F := \lambda x.\lambda y.y$ (заметим, что $T = K$);
\item $Not := \lambda x.x\ F\ T$;
\item $1 := \lambda f.\lambda x.f\ x$;
\item $(+1) := \lambda n.\lambda f.\lambda x.n\ f\ (f\ x)$.
\end{enumerate}

\item Рассмотрим альтернативный комбинаторный базис $BCKW$, три новых комбинатора понимаются нами так:
$$\begin{array}{ll}
B := \lambda x.\lambda y.\lambda z.x\ (y\ z) & \text{Шейнфинкель: Z, Zusammensetzungsfunkion}\\
C := \lambda x.\lambda y.\lambda z.x\ z\ y & \text{Шейнфинкель: T, Vertauschungsfunkion}\\
W := \lambda x.\lambda y.x\ y\ y
\end{array}$$

\begin{itemize}
\item Выразите $I$ в этом базисе;
\item Выразите $F$ в этом базисе;
\item Выразите $Not$ в этом базисе;
\item Покажите, что данный базис позволяет выразить любой комбинатор из лямбда-исчисления;
\item Найдите схемы аксиом, соответствующие данным комбинаторам в исчислении в Гильбертовском стиле;
также покажите, что импликативный фрагамент ИИВ, в котором схемы аксиом заменены на найденные вами,
корректен и полон.
\end{itemize}

%\item Как известно, нормальный порядок редукции медленнее аппликативного. Чтобы (частично) преодолеть
%эту проблему, давайте рассмотрим правило \emph{мемоизации}:
%
%При редукции терма $(\lambda x.P)Q$ не производится текстовой подстановки терма $Q$ вместо всех вхождений $x$ в $P$.
%Вместо этого мы запоминаем ссылку на выражение $P$ вместо всех вхождений $x$. Такая ссылка <<прозрачна>> на чтение,
%но если в ходе бета-редукции подставленного выражения редукция происходит внутри выражения по ссылке, редукция
%касается всех вхождений сразу.
%
%Например: $(\lambda x.x x x) (I I) \rightarrow_\beta \xi\ \xi\ \xi [\xi := I I] \rightarrow_\beta \xi\ \xi \xi [\xi := I]$
%
%Если же в ходе редукции замена пересекает границу ссылки --- 

\item Найдите \emph{замкнутое} выражение, в котором при редукции какого-то подвыражения вида $(\lambda.P)\ Q$ (при 
подстановке $Q$ вместо переменной $0$ в $P$) требуется изменение индексов в $Q$ в нотации де Брауна.
Предложите алгоритм для этого и покажите его корректность.
Также, предложите способ, как в нотации де Брауна указывать глобальные переменные.

\item Покажите, что расширенные полиномы (см. теорему о выразительной силе просто-типизированного лямбда исчисления)
выражаются в данном исчислении (т.е. для каждого E(m,n) найдётся лямбда-выражение $E: \eta\rightarrow\eta\rightarrow\eta$,
что $\overline{E(m,n)} = E\ \overline{m}\ \overline{n}$). 
В частности, важным результатом в этом задании является лямбда-терм для умножения:
$(\cdot): \eta\rightarrow\eta$, что $(\cdot)\ \overline{m}\ \overline{n} =_\beta \overline{m\cdot n}$.
Также найдите другой терм для умножения, который не имеет типа в просто-типизированном лямбда исчислении.

\item Пусть $\neg\alpha := \alpha\rightarrow\bot$. Также пусть $\varphi\leftrightarrow\psi := (\varphi\rightarrow\psi)\with(\psi\rightarrow\varphi)$.
Найдите лямбда-терм в полном просто-типизированном лямбда исчислении, обитающий в типе:
\begin{itemize}
\item $(\alpha\vee\beta)\rightarrow\neg(\neg\alpha\with\neg\beta)$ --- правило де Моргана;
\item $(\alpha\rightarrow\beta\rightarrow\gamma) \leftrightarrow (\alpha\with\beta\rightarrow\gamma)$ --- карринг;
\item $(\alpha\with(\beta\vee\gamma)) \leftrightarrow ((\alpha\vee\beta)\with(\alpha\vee\gamma))$ --- дистрибутивность.
\end{itemize}

\item Обитаем ли тип $(\alpha\rightarrow\beta)\rightarrow((\alpha\rightarrow\bot)\vee\beta)$? Ответ докажите.
\end{enumerate}

\section*{Задание №3. Система F.}
\begin{enumerate}
\item На лекции алгоритм реконструкции типов для просто типизированного лямбда-исчисления формулировался для исчисления по Карри.
Сформулируйте алгоритм реконструкции типов для исчисления по Чёрчу --- чем он будет отличаться от алгоритма для исчисления по Карри?
\item Постройте систему уравнений в алгебраических термах для $Y$-комбинатора, покажите
её несовместность.
\item Задача, оставшаяся с практики 21.03.2026: в каком соотношении находятся 
$\llbracket \alpha\rightarrow\beta\rrbracket$ и $\llbracket \alpha\rightarrow\alpha \rrbracket$ ---
есть ли вложенность одного множества в другое?

\item Рассмотрим более общий тип чёрчевского нумерала для системы F:
$\forall \alpha.(\alpha\rightarrow\alpha)\rightarrow(\alpha\rightarrow\alpha)$.
Постройте значения этого нумерала, постройте операции сложения, умножения и возведения в 
степень для него, докажите их тип в системе F.

\item Выразите следующие связки через $\rightarrow, \forall$, реализуйте соответствующие функции, докажите их типы в $F$:
\begin{enumerate}
\item $\with$, реализуйте функции MkPair, $\pi_L$, $\pi_R$;
\item $\vee$, реализуйте функции $in_L$, $in_R$, Case.
\end{enumerate}
\item Рассмотрите функции. 
Каждая из функций использует конъюнкцию или дизъюнкцию, их можно выразить через квартор всеобщности и импликацию.
Однако, получившееся выражение может не соответствовать системе HM. В связи с этим:
реализуйте их в системе F и примените алгоритм W (получится ли это?); определите ранг типа этих функций;
реализуйте эти функции на Хаскеле.
Пользуйтесь аннотацией \verb!{-# LANGUAGE RankNTypes #-}!, позволяющей обойти ограничения HM (за счёт потери
разрешимости типовой системы языка), смотрите лекцию 5.
\begin{enumerate}
\item функция $\pi_L: \alpha\with\beta\rightarrow\alpha$ (левая проекция для упорядоченной пары) --- раскройте её с учётом способа
выражения конъюнкции через квантор всеобщности и импликацию.
\item функция $in_L: \alpha\rightarrow\alpha\vee\beta$
\item $\texttt{let }id = \lambda x.x\texttt{ in }\langle \texttt{id }0, \texttt{id ``a''}\rangle$
\end{enumerate}

\item Модифицируйте алгоритм W, чтобы включить в него работу с упорядоченными парами и алгебраическими типами.

\item Запишите тип списка (из целых чисел, используйте константу \texttt{int}) в HM, а также реализуйте функции для получения его длины и функцию map.
Используйте расширения (типизацию для Y-комбинатора, мю-оператор в варианте изорекурсивных типов).

\item Выразите Y комбинатор в Хаскеле (Окамле), укажите его тип.

%\item Покажите, что если $rk(\sigma,1)$, то для выражения $\sigma$ найдётся эквивалентное $\sigma'$ с поверхностными
%кванторами. 

%\item Покажите, что в предыдущем задании также имеется изоморфизм типов: существует биективная функция $\sigma\rightarrow\sigma'$,
%которую можно выразить лямбда-выражениями.

\item Рассмотрим QuickSort:
\begin{verbatim}
let rec quick l =  match l with
    [] -> []
  | l1 :: ls -> List.filter (fun x -> x < l1) ls @ [l1] @ List.filter (fun x -> x >= l1)
\end{verbatim}

Укажите полные типовые схемы в системе HM для всех функций, участвующих в данном примере (тип списка раскрывать не надо).

\item Заметим, что список --- это <<параметризованные>> числа в 
аксиоматике Пеано. Число --- это длина списка, а к каждому штриху мы присоединяем какое-то значение.
Операции добавления и удаления элемента из списка --- это операции прибавления и вычитания
единицы к числу.

Рассмотрим тип <<бинарного списка>>:

\begin{verbatim}
type 'a bin_list = Nil | Zero of (('a*'a) bin_list) | One of 'a * (('a*'a) bin_list);;
\end{verbatim}

Заметим, что здесь мы рассматриваем двоичную запись числа (чередующиеся \verb!Zero! и \verb!One!) ---
двоичную запись длины бинарного списка, и элемент двоичной записи номер $n$ хранит $2^n$ или $2^n+1$ 
значение (в зависимости от типа элемента). Например, 5-элементный список:

\begin{verbatim}
One ("a", Zero (One ((("b","c"),("d","e")), Nil)))
\end{verbatim}

По идее, операция добавления элемента к списку записывается на языке Окамль вот так 
(сравните с прибавлением 1 к числу в двоичной системе счисления):

\begin{verbatim}
let rec add elem lst = match lst with
    Nil -> One (elem,Nil)
  | Zero tl -> One (elem,tl)
  | One (hd,tl) -> Zero (add (elem,hd) tl)
\end{verbatim}

Однако, тип этой функции Окамль вывести автоматически не может, его надо указывать явно,
чтобы код скомпилировался:

\begin{verbatim}
let rec add : 'a . 'a -> 'a bin_list -> 'a bin_list = fun elem lst -> match lst with
\end{verbatim}
\begin{enumerate}
\item Какой тип имеет \verb!add! в (расширенной) системе $F$ (напомним, поскольку функция рекурсивна,
она должна использовать Y-комбинатор в своём определении)?
Считайте, что семейство типов \verb!bin_list 'a! предопределено и обозначается как $\tau_\alpha$.
Также считайте, что определены функции roll и unroll с надлежащими типами.
Какой ранг имеет тип этой функции? Почему этот тип не выразим в типовой системе Хиндли-Милнера?
\item Предложите функции для печати списка и для удаления элемента списка (головы).
\item Предложите функцию для эффективного соединения двух списков (источник для 
вдохновения --- сложение двух чисел в столбик).
\item Предложите функцию для эффективного выделения $n$-го элемента из списка.
\end{enumerate}


\item Предложите выражение на языке C++ (возможно, использующее шаблоны), имеющее следующий род (тип):
\begin{enumerate}
\item $\star\rightarrow\star\rightarrow\star$; $\ \star\rightarrow\textbf{unsigned}$
\item $\textbf{int}\rightarrow(\star\rightarrow\star)$
\item $(\star\rightarrow\textbf{int})\rightarrow\star$
\item $\Pi x^\star.\lambda n^\textbf{int}.F(n,x)$, где $$F(n,x) = \left\{\begin{array}{ll}\textbf{int}, & n = 0\\
                                   x\rightarrow F(n-1,x), & n > 0\end{array}\right.$$
\end{enumerate}

\end{enumerate}

\section*{Задание №4. Язык Аренд.}
\begin{enumerate}
\item Докажите, приведя компилирующуюся программу на языке Аренд (возможно, вам потребуются функции и приёмы, 
изложенные в документации по языку: \url{https://arend-lang.github.io/documentation/tutorial/PartI/}):
\begin{enumerate}
%\item ассоциативность сложения;
%\item коммутативность сложения;
\item коммутативность умножения;
\item дистрибутивность: $(a + b)\cdot c = a\cdot c + b \cdot c$;
\item куб суммы: $(a + 1)^3 = a^3 + 3\cdot a^2 + 3 \cdot a + 1$.
\end{enumerate}

\item Определим, что $x$ делится на $p$, если обитаем тип \verb!\Sigma (q : Nat) (p * q = x)!.
\begin{enumerate}
\item Покажите, что если $x$ делится на 6, то $x$ делится и на 3;
\item Покажите, что $x!$ делится на $x$;
\item Покажите, что если $x$ делится на $y$ и $y$ делится на $z$, то $x$ делится на $z$;
\end{enumerate}

\item Определите предикат (т.е. функцию с надлежащим типом) для формализации понятия простого числа \verb!isPrime!.
Покажите, что:
\begin{enumerate}
\item 3 и 11 --- простые числа;
\item Произведение простых чисел непросто;
\item 2 --- единственное чётное простое число.
\end{enumerate}

\item Определим отношение <<меньше>> на натуральных числах так (с помощью индуктивного типа, обобщения алгебраического типа данных ---
зависимого типа, в котором при разных значениях аргументов типа допустимы разные конструкторы):
\begin{verbatim}
\data NatLessEq (a b : Nat) \with
  | 0, m => natlesseq-zero
  | suc m, suc n => natlesseq-next (NatLessEq m n)
\end{verbatim}

Например, конструктор \verb!natlesseq-zero! можно использовать только если первый аргумент типа --- число 0.
А конструктор \verb!natlesseq-next! применим только если первый аргумент больше 0; при этом данный конструктор
требует значение типа с определёнными аргументами в качестве своего аргумента.

Будем говорить, что $a \preceq b$ тогда и только тогда, когда \verb!NatLessEq a b! обитаем.
Например, утверждение $1 \preceq 3$ доказывается так:

\begin{center}\verb!\func one-le-three : NatLessEq 1 3 => natlesseq-next (natlesseq-zero)!\end{center}

\emph{В самом деле}, \verb!natlesseq-zero! может являться конструктором типа \verb!NatLessEq 0 b!,
а тогда \begin{center}\verb!natlesseq-next (natlesseq-zero) : NatLessEq 1 (b+1)!\end{center} Унифицировать $b+1$ и $3$
компилятор (в данном случае) может самостоятельно, и потому код выше проходит проверку на корректность.

Докажите (везде предполагается, что $a,b,c : \texttt{Nat}$, если не указано иного):
\begin{enumerate}
\item $a \preceq b$ тогда и только тогда, когда $a$ меньше или равно $b$ в смысле натуральных чисел (здесь требуется рассуждение
на мета-языке).
\item $a \preceq a + b + 1$; то есть, определите функцию\\\verb!\func n-less-sum (a b : Nat) : NatLessEq a (a Nat.+ suc b)!
\item Если $a \preceq b$, то $a + c \preceq b + c$
\item Если $a \preceq b$ и $c \preceq d$, то $a \cdot c \preceq b \cdot d$
\item $a \preceq 2^a$
\item Транзитивность: если $a \preceq b$ и $b \preceq c$, то $a \preceq c$
\item $a \preceq b \vee b \preceq a$
\item Найдите стандартное определение отношения <<меньше>> в библиотеке Аренда (\verb!Nat.<!) и докажите, что $a \preceq b$ 
тогда и только тогда, когда $a < b$ или $a = b$ (реализуйте функции \\\verb!there (p : NatLessEq a b) : Data.Or (a Nat.< b) (a = b)!
и обратную к ней).
\item Покажите, что $a \preceq b$ тогда и только тогда, когда $\exists k^{\mathbb{N}_0}.a + k = b$.
\end{enumerate}

\item С точки зрения изоморфизма Карри-Ховарда индуктивные типы можно воспринимать как аналоги предикатов.
В этом задании надо построить индуктивные типы для различных предикатов:
\begin{enumerate}
\item Факториал (\verb!IsFact n!), который будет обитаем только для таких $n$, что $n = k!$.
Докажите на языке Аренд, что \verb!IsFact! $(1 \cdot 2 \cdot 3 \cdot \dots \cdot n)$ всегда обитаем, а
тип \verb!IsFact 3! --- необитаем.
\item Наибольший общий делитель двух чисел \verb!GCD x a b!; \emph{указание/пожелание:} воспользуйтесь алгоритмом Эвклида.  
\item Ограниченное натуральное число \verb!IndFin n!, обитателями типа являются только те числа, которые меньше $n$.
В стандартной библиотеке \verb!Fin! определён через натуральные числа; 
сделайте это исключительно через индуктивные типы. Покажите, что если \verb!x : IndFin m! и \verb!y : IndFin n!,
то \verb!x + y : IndFin (m + n)!.
\end{enumerate}

\item Рассмотрим следующее доказательство уникальности элементов списка:
\begin{verbatim}
\func not-member (A : \Type) (elem : A) (l : List A) : \Type \elim l
  | nil => \Sigma
  | :: hd tl => \Sigma (Not (hd = elem)) (not-member A elem tl)

\func unique-list (A : \Type) (l : List A) : \Type \elim l
  | nil => \Sigma
  | :: hd tl => \Sigma (not-member A hd tl) (unique-list A tl)
\end{verbatim}

Докажем, что список $[0,1,2]$ состоит из уникальных элементов:
\begin{verbatim}
\func r-unique => unique-list Nat (0 :: 1 :: 2 :: nil)
\func x : r-unique => ((contradiction, (contradiction, ())), ((contradiction, ()), ((), ())))
\end{verbatim}

\begin{enumerate} 
\item Напишите функцию, строящую список \verb!b! натуральных чисел от 0 до n и доказательство\\ \verb!unique-list Nat b!.
\item Покажите, что если $a_0 < a_1 < \dots < a_n$, то список $[a_0,a_1,\dots,a_n]$ уникальный.
\item Покажите, что если $n \ge 2$ и $a_k \ne a_{k+1}$ при $0 \le k < n$, то необязательно, что список
$[a_0,a_1,\dots,a_n]$ --- уникальный.
\end{enumerate}

\item Рассмотрим правый дальний нижний угол $\lambda$-куба ($\{(\star,\star);(\star,\square);(\square,\star)\}$).
Можно предположить, что тогда в такой системе возможны и функции рода $f: \star\rightarrow\star$ 
(как композиция функций $p: \star\rightarrow\alpha$ и $q: \alpha\rightarrow\star$ ---
например, можно кодировать тип его именем, затем по имени типа восстанавливать сам тип обратно).
Почему всё-таки такие функции в обобщённых типовых системах невозможны без четвёртого элемента $(\square,\square)$?

\item Какова должна быть топология на множестве
пар натуральных чисел (интуитивно мы будем понимать эти пары как рациональные числа, пары <<числитель-знаменатель>>), 
чтобы непрерывными были бы те и только те функции, для которых выполнено $f(p,q) = f(p',q')$ для всех 
таких $p,p',q$ и $q'$, что $p\cdot q' = p'\cdot q$. Напомним, что равенство мы понимаем как наличие непрерывного пути 
между точками.
\end{enumerate}

\section*{Задание №5. Дополнительные задачи на доказательства. Равенство.}

\begin{enumerate}

\item Рассмотрим модуль числа, $|x|$.
\begin{enumerate}
\item Предложите определение этой конструкции (должна отображать \verb!Int! в \verb!Nat!). Докажите, что $||x|| = |x|$.
\item Докажите, что если $x \ge 0$, то $|x|=x$.
\item Докажите, что $|a - b| + |a + b| = 2\cdot\max(|a|,|b|)$.
\end{enumerate}

\item На практике было много вопросов по поводу доказательства 
$\forall x\in\mathbb{N}.\forall y\in\mathbb{N}.x < y \rightarrow \neg\exists z\in\mathbb{N}.x = z\cdot y$.
Рассмотрим несколько упражнений (дополнительных лемм), которые могут помочь в этом доказательстве.
\begin{enumerate}
\item Заметим, что это утверждение неверно при $x,y,z\in\mathbb{N}_0$.
Поэтому, предложите контрпример и докажите его на языке Аренд.
\item Рассмотрим деление с остатком, divmod:\\ 
\verb!\divmod (a b : Nat) (b /= 0) => \Sigma (q r : Nat) (a = q * b + r) (r < b)!.
Будем обозначать первый элемент пары (собственно деление) как $//$ (в стиле языка Питон).
Докажите, что $a_1 // b \geq a_2 // b$, если $a_1 > a_2$.
\item Докажите, что $a_1 // b > a_2 \div b$, если $a_1 \geq a_2 + b$.
\item Докажите, что $a // b_1 \geq a \div b_2$, если $b_1 < b_2$.
\item Докажите, что если $a_1 // b = a_2 // b$, то $|a_2 - a_1| < b$.
\item Докажите, что если $(p1,q1,r1) = \texttt{divmod}\ a\ b$ и $(p2,q2,r2) = \texttt{divmod}\ a\ b$, то $p1=p2$ и $q1=q2$.
\item Докажите, что если $a // b = 1$, то $a < 2b$ (ну, или $a = 1\cdot b + k, k < b$).
\item Докажите, что если $a // b = 0$, то $a < b$ и $a\ \texttt{mod}\ b = a$.
\item Докажите, что если $a < b$, то $a\ \texttt{mod}\ b = a$ и $a // b = 0$.
\end{enumerate}

\item Как вы помните из лекции, в языке Аренд существует иерархия вложенных типовых универсумов. Если в типе отсутствует
упоминание \verb!\Type!, то данный тип принадлежит универсуму 0. Однако, если в типе упоминается \verb!\Type k!,
то тип принадлежит универсумам, не меньшим $k+1$. Уровень универума обозначается специальным ключевым словом \verb!\lp!.
Над индексами можно проводить простые операции и сопоставление с образцом (\verb!\suc\lp!). Более подробно
можно это прочесть в документации по языку Аренд:

\url{https://arend-lang.github.io/documentation/tutorial/PartI/universes.html}
%\begin{enumerate}
%\item Определите функцию \verb!\id!, возвращающую 
%\end{enumerate}

Рассмотрим определения:
\begin{verbatim}
\func ChurchT (x : \Type) => (x -> x) -> (x -> x)
\func Church => \Pi (x : \Type) -> ChurchT x
\func Zero : Church => \lam t f x => x

\func incT (t : \Type) (n : ChurchT t) => \lam f x => n f (f x)
\func pair_plus (type : \Type) (pair : \Sigma (ChurchT type) (ChurchT type)) :
  \Sigma (ChurchT type) (ChurchT type) => (pair.2, incT type pair.2)
\func dec (n : Church) : Church => \lam (t : \Type) =>
    (n (\Sigma (ChurchT t) (ChurchT t)) (pair_plus t) (Zero t, Zero t)).1
\func sub (a : Church (\suc\lp)) (b : Church) => a (Church \lp) dec b
\end{verbatim}

Определите, развивая определения выше:
\begin{enumerate}
\item Операцию умножения.
\item Операцию <<деление на три>> (естественно, в версии, не использующей $Y$-комбинатор).
\item Операцию возведения в степень, определявшуюся как $\lambda m.\lambda n.n\ m$.
\item Деление.
\item Вычисление факториала.
\end{enumerate}

\item Введём тип данных \verb!IsEven!:

\begin{verbatim}
\data IsEven (n : Nat) \elim n
  | 0 => zero-is-even
  | (suc (suc k)) => next-next (IsEven k)
\end{verbatim}

Доказать, что этот тип является утверждением, можно например так:

\begin{verbatim}
\func all-even-different (n : Nat) (a b : IsEven n) : a = b \elim n, a, b
  | 0, zero-is-even, zero-is-even => idp
  | suc (suc n), next-next a, next-next b => pmap next-next (all-even-different n a b)

\func is-even-isProp (n : Nat) : isProp (IsEven n) => all-even-different n
\end{verbatim}

Обратите внимание: по переменным $n$, $a$ и $b$ производится \emph{элиминация}, то есть множество 
значений переменных разбивается на фрагменты (в соответствии с конструкторами типа данных), 
и доказательство утверждения проводится независимо для каждого фрагмента; 
в частности, цель доказательства изменяется и просходит замена переменных $a$ и $b$ на сопоставляемые
варианты (вместо ожидаемого типа \verb!a = b! мы ожидаем тип \verb!zero-is-even = zero-is-even!, и т.п.).
Сравните с лекцией про элиминаторы, каждый из вариантов --- тело отдельной функции-элиминатора. 

Чтобы воспроизвести тот же эффект в конструкции \verb!\case!, нужно указывать ключевое слово 
\verb!\elim! перед каждой элиминируемой переменной: 
\begin{verbatim}
\case \elim n, \elim a, \elim b \with { ... }
\end{verbatim}

Полный код, определяющий тип \verb!IsEven! (вместе с доказательством того, что тип --- утверждение),
выглядит так:

\begin{verbatim}
\data IsEven (n : Nat) \elim n
  | 0 => zero-is-even
  | (suc (suc k)) => next-next (IsEven k)
  \where {
    \func all-even-different ... -- скопируйте код функции сюда
    \use \level is-even-isProp (n : Nat) : isProp (IsEven n) => all-even-different n
  }
\end{verbatim}

Однако, незавершённым остаётся доказательство разрешимости типа \verb!IsEven n!. Восполните лакуны:

\begin{verbatim}
\func even-is-dec (a : Nat) : Dec (IsEven a) \elim a
  | 0 => yes zero-is-even
  | 1 => no {?}
  | suc (suc a) => {?}
\end{verbatim}

\item Справедливости ради, в предыдущем задании компилятор сам может догадаться, что \verb!IsEven! --- утверждение.
Но далеко не для всех типов это очевидно, и тогда функция с префиксом \verb!\use \level! становится необходимой.
Например, в следующем типе <<простое число>> данная функция должна доказать, что любые два значения типа при данном $x$ равны:

\begin{verbatim}
\data Div3 (x : Nat)
  | remainder-zero (Exists (p : Nat) (p Nat.* 3 = x))
  | remainder-one (Exists (p : Nat) (p Nat.* 3 Nat.+ 1 = x))
  | remainder-two (Exists (p : Nat) (p Nat.* 3 Nat.+ 2 = x))
  \where \use \level levelProp {x : Nat} (a b : Div3 x) : a = b => {?}
\end{verbatim}

\begin{enumerate}
\item Замените \verb!{?}! в тексте выше на корректное доказательство.
\item Определите функцию, которая бы по $x$ и значению $\exists p q.3 \cdot p + q = x \with 0 \le q < 3$ возвращала бы \verb!Div3 x!
(понятно, можно разделить $x$ на 3, но нам уже результат деления дали вторым аргументом --- задача в том, чтобы им воспользоваться).
\item Постройте аналогичный тип \verb!Prime x! для простых чисел --- с тремя вариантами \verb!less-than-two!, \verb!is-prime!, 
\verb!is-composite! --- и определите функцию, строящую по числу значение данного типа.
\item Покажите, что в типе \verb!Prime (x*x + 2*x + 1)! всегда (кроме граничных случаев) обитает вариант \verb!is-composite!.
\item Покажите, что в типе 
\begin{verbatim}
\data SuperDec (P : \Prop)
| sure P
| nope (P -> Empty)
| neither ((P || (P -> Empty)) -> Empty)
\where \use \level superDecProp {P : \Prop} (a b : SuperDec P): a = b => {?}
\end{verbatim}
вариант \verb!neither! невозможен (также, заполните пропущенное доказательство \verb!superDecProp!).
\end{enumerate}
\end{enumerate}

\end{document}
